\section{Logical Design}

\subsection{Volume Table}
Volumes are estimated for one year of operation.

\begin{table}[H]
    \scriptsize
    \renewcommand{\arraystretch}{1.3}
    \begin{tabularx}{\textwidth}{|X|X|X|X|}
    \hline
    \textbf{Concept} & \textbf{Type} & \textbf{Volume} & \textbf{Notes} \\ \hline
    Customer        & E & 10.000  & Assumption \\ \hline
    Event\_Location & E & 8.000   & $\sim$0.8 locations per customer on average \\ \hline
    Booking         & E & 50.000  & 10 bookings/day $\times$ 365 days $\times$ overheads \\ \hline
    Team            & E & 20      & Assumption \\ \hline
    Member          & E & 100     & 5 members per team $\times$ 20 teams \\ \hline
    Depot           & E & 10      & Assumption \\ \hline
    Municipality    & E & 30      & $\sim$3 municipalities per depot \\ \hline
    PLACE           & R & 8.000   & Customer $\to$ Event\_Location (1:N) \\ \hline
    RESERVATION     & R & 50.000  & Customer $\to$ Booking (1:N) \\ \hline
    BL              & R & 50.000  & Event\_Location $\to$ Booking (1:N) \\ \hline
    HANDLED         & R & 50.000  & Booking $\to$ Team (N:1) \\ \hline
    HAS             & R & 100     & Team $\to$ Member (1:N) \\ \hline
    BELONG          & R & 40      & $\sim$2 depots per team on average (N:M) \\ \hline
    COVER           & R & 30      & Depot $\to$ Municipality (1:N) \\ \hline
    \end{tabularx}
    \caption{Volume Table}
\end{table}

\subsection{Access Table}

\noindent\textbf{Operation 1:} Register a new customer (30 times per day)
\begin{table}[H]
    \renewcommand{\arraystretch}{1.3}
    \begin{tabularx}{\textwidth}{|X|X|X|X|}
    \hline
    \textbf{Concept} & \textbf{Type} & \textbf{Accesses} & \textbf{R/W} \\ \hline
    Customer & E & 1 & W \\ \hline
    \end{tabularx}
    \caption{Access Table -- Operation 1}
\end{table}
\noindent Total Cost: $2 \times 30 = 60$ per day\\

\noindent\textbf{Operation 2:} Add a new booking (10 times per day)
\begin{table}[H]
    \renewcommand{\arraystretch}{1.3}
    \begin{tabularx}{\textwidth}{|X|X|X|X|}
    \hline
    \textbf{Concept} & \textbf{Type} & \textbf{Accesses} & \textbf{R/W} \\ \hline
    Booking         & E & 1    & W \\ \hline
    Team            & E & 1    & R \\ \hline
    Customer        & E & 1    & R \\ \hline
    Event\_Location & E & 1    & R \\ \hline
    Team            & E & 1    & W (Installation\_Number update) \\ \hline
    \end{tabularx}
    \caption{Access Table -- Operation 2}
\end{table}
\noindent Total Cost: $(2 + 1 + 1 + 1 + 2) \times 10 = 70$ per day\\

\noindent\textbf{Operation 3:} Add a new team (5 times per day)
\begin{table}[H]
    \renewcommand{\arraystretch}{1.3}
    \begin{tabularx}{\textwidth}{|X|X|X|X|}
    \hline
    \textbf{Concept} & \textbf{Type} & \textbf{Accesses} & \textbf{R/W} \\ \hline
    Team & E & 1 & W \\ \hline
    \end{tabularx}
    \caption{Access Table -- Operation 3}
\end{table}
\noindent Total Cost: $2 \times 5 = 10$ per day\\

\noindent\textbf{Operation 4:} Add a new member to a team (2 times per day)
\begin{table}[H]
    \renewcommand{\arraystretch}{1.3}
    \begin{tabularx}{\textwidth}{|X|X|X|X|}
    \hline
    \textbf{Concept} & \textbf{Type} & \textbf{Accesses} & \textbf{R/W} \\ \hline
    Team   & E & 1    & R (existence check) \\ \hline
    Member & E & 5    & R (capacity check: count members in team) \\ \hline
    Member & E & 1    & W \\ \hline
    \end{tabularx}
    \caption{Access Table -- Operation 4}
\end{table}
\noindent Total Cost: $(1 + 5 + 2) \times 2 = 16$ per day\\

\noindent\textbf{Operation 5:} Get all members of a specific team (5 times per day) $\star$
\begin{table}[H]
    \renewcommand{\arraystretch}{1.3}
    \begin{tabularx}{\textwidth}{|X|X|X|X|}
    \hline
    \textbf{Concept} & \textbf{Type} & \textbf{Accesses} & \textbf{R/W} \\ \hline
    Member & E & 5 & R (avg 5 members per team) \\ \hline
    \end{tabularx}
    \caption{Access Table -- Operation 5 (most frequent query)}
\end{table}
\noindent Total Cost: $5 \times 5 = 25$ per day\\

\subsubsection{Redundancy Analysis: \texttt{Installation\_Number}}

\noindent The attribute \texttt{Installation\_Number} in \textbf{Team} stores the number of bookings handled
by that team. This is a \textbf{derived, redundant} attribute. We analyse whether to keep it.\\

\noindent\textit{Without redundancy} -- to compute Installation\_Number for a team we must scan Booking:
\begin{table}[H]
    \renewcommand{\arraystretch}{1.3}
    \begin{tabularx}{\textwidth}{|X|X|X|X|}
    \hline
    \textbf{Concept} & \textbf{Type} & \textbf{Accesses} & \textbf{R/W} \\ \hline
    Booking & E & 2.500 & R (50.000 bookings / 20 teams) \\ \hline
    \end{tabularx}
    \caption{Cost of computing Installation\_Number without redundancy}
\end{table}

\noindent\textit{With redundancy} -- Operation 2 must update Team in addition to inserting Booking:
\begin{table}[H]
    \renewcommand{\arraystretch}{1.3}
    \begin{tabularx}{\textwidth}{|X|X|X|X|}
    \hline
    \textbf{Concept} & \textbf{Type} & \textbf{Accesses} & \textbf{R/W} \\ \hline
    Team & E & 1 & W (extra write per booking insert) \\ \hline
    \end{tabularx}
    \caption{Extra cost of maintaining Installation\_Number (redundancy)}
\end{table}

\noindent Assuming Installation\_Number is queried at most $q$ times per day, the break-even point is:
\[
q \times 2.500 > 10 \times 2 \quad \Rightarrow \quad q > 0.008
\]
Since $q \geq 1$, it is \textbf{always convenient to keep the redundancy}. The extra write cost of the
trigger is negligible ($+2 \times 10 = 20$ accesses/day) compared to the savings on read queries.
The redundancy is maintained automatically by \texttt{trg\_update\_installation\_count}.


\subsection{UML Schema}

\begin{figure}[H]
    \centering
    \includegraphics[width=\textwidth]{images/UML.png}
    \caption{UML Schema -- Logical design of SageFit}
\end{figure}

\subsection{Relational Schema}

\begin{table}[H]
    \renewcommand{\arraystretch}{1.3}
    \begin{tabularx}{\textwidth}{|X|}
    \hline
    \texttt{Customer(\underline{ID}, Name, Type, Email, Phone, Address)} \\ \hline
    \texttt{Event\_Location(\underline{ID}, Street, House\_Number, Postal\_Code, City, Province,} \\
    \quad \texttt{Setup\_Time\_Estimate, Equipment\_Capacity, Customer\_ID)} \\ \hline
    \texttt{Booking(\underline{ID}, Type, Booking\_Date, Duration, Cost, Booking\_Method,} \\
    \quad \texttt{Contract\_Type, Team\_Code, Customer\_ID, Location\_ID)} \\ \hline
    \texttt{Team(\underline{Code}, Name, Max\_Members, Installation\_Number)} \\ \hline
    \texttt{Member(\underline{ID}, Name, Surname, Phone, Email, Team\_Code)} \\ \hline
    \texttt{Depot(\underline{ID}, Name, City, Province, Address, Region, Number\_Of\_Employees)} \\ \hline
    \texttt{Municipality(\underline{ID}, Name, Depot\_ID)} \\ \hline
    \texttt{Team\_Depot(\underline{Team\_Code, Depot\_ID})} \\ \hline
    \end{tabularx}
    \caption{Relational Schema (underlined attributes = primary key)}
\end{table}

\subsubsection{Translation Rules Applied}
\begin{enumerate}
    \item \textbf{Customer $\to$ Event\_Location (1:N PLACE)}: FK \texttt{Customer\_ID} added to \texttt{Event\_Location}
    \item \textbf{Customer $\to$ Booking (1:N RESERVATION)}: FK \texttt{Customer\_ID} added to \texttt{Booking}
    \item \textbf{Event\_Location $\to$ Booking (1:N BL)}: FK \texttt{Location\_ID} added to \texttt{Booking}
    \item \textbf{Team $\to$ Booking (1:N HANDLED)}: FK \texttt{Team\_Code} added to \texttt{Booking}
    \item \textbf{Team $\to$ Member (1:N HAS)}: FK \texttt{Team\_Code} added to \texttt{Member}
    \item \textbf{Team $\leftrightarrow$ Depot (N:M BELONG)}: New join table \texttt{Team\_Depot} with both FKs
    \item \textbf{Depot $\to$ Municipality (1:N COVER)}: FK \texttt{Depot\_ID} added to \texttt{Municipality}
\end{enumerate}
